%%% SRS - Organisation: Beginn
\begin{center}
  \makeatletter
  \textbf{
    \@title
  }
  \makeatother
  \\[1em]
  Informatik -- Eingebettete Systeme -- Hochschule Bremerhaven
\end{center}

\begin{table}[ht]
  \centering
  \begin{tabular}[t]{|p{15mm}|p{30mm}|p{30mm}|p{50mm}|r|}
    \hline
    \textbf{Version} & \textbf{Autor*inn*en} & \textbf{geändert} & \textbf{Beschreibung} & \textbf{fertig am}\\
    \hline
    0.1              & A. B.            &  Abschnitte 5 bis 7  & Erstellung des Dokumentes & 01.10.2019\\
    \hline
  \end{tabular}
  \caption{Versionshistorie}
\end{table}

\vspace*{4em}

Relevante Standards und Normen
\begin{itemize}
\item IEEE 29148-2018 - ISO/IEC/IEEE International Standard - Systems
  and software engineering -- Life cycle processes -- Requirements
  engineering

  \cite{ieee-29148-2018-requirements-engineering}
\item ISO-9126 / IEC-25012 -- Software- und Systemqualitätsmerkmale

  \cite{iso-iec-25012-2008-product-quality}
\end{itemize}
\clearpage
%%% SRS - Organisation: Ende

\chapter{Einleitung}
\label{srs:einleitung}

\hinweis{Dieser Text ist nur für Sie. Am Ende des Semesters muss
  dieser Text hier entfernt sein. Die System-Anforderungs\-spezifikation
  bündelt alle Anforderungen an das fertige Produkt und alle
  Projekterzeugnisse. Anforderungen kann man gewöhnlich nicht einfach
  auflisten. Viele Anforderungen widersprechen einander und so müssen
  sie priorisiert werden. Dies geschieht über Qualitätsszenarien,
  einfache Sätze, die beschreiben, wer etwas tut, womit er oder sie
  das tut und, was dabei herauskommen soll. Letzteres wird in einer
  messbaren Art und Weise beschrieben.}

\section{Projektübersicht}
\label{srs:projektuebersicht}

\hinweis{Diese Übersicht geht auf die Anforderungen ein.}

\begin{figure}[htp]
  \centering

  \textbf{\hinweis{Hierhin gehört der Systemkontext, der alle Bestandteile
    zeigt, die über Schnittstellen mit dem System interagieren.}}

  \caption{Systemkontext}
  \label{srs:systemkontext}
\end{figure}

\subsection{Qualtitätsszenarien}

\section{Produktumfang und -einsatz}
\hinweis{Hierhin gehören Qualitäten des Produkts. Diese müssen in den
  Anforderungen referenziert werden, um die Anforderungen zu
  begründen.}

\chapter{Anforderungen}
\label{srs:anforderungen}
\hinweis{Die Anforderungen sind durchnummeriert und beginnen alle mit dem
Präfix \textbf{REQ}.}

\section{Anforderungen externer Schnittstellen}
\label{srs:anforderungen-externer-schnittstellen}

\hinweis{(Präfix: REQ-EXT)}

\subsection{Benutzerschnittstellen}
\label{srs:benutzerschnittstellen}

\hinweis{(Präfix: REQ-EXT-USER)}
\subsubsection{REQ-EXT-USER 1} \hinweis{Benutzer muss mit dem System
  durch Tasterdruck und Rückmeldung durch eine Anzeige interagieren
  können.}

\paragraph{Betroffene Qualitätszenarien} \hinweis{QS-1, QS-2, QS-5}

\subsection{Hardwareschnittstellen}
\label{srs:hardwareschnittstellen}

\hinweis{(Präfix: REQ-EXT-HW)}

\subsection{Softwareschnittstellen}
\label{srs:softwareschnittstellen}

\hinweis{(Präfix: REQ-EXT-SW)}

\subsection{Kommunikationsprotokolle}
\label{srs:kommunikationsprotokolle}

\hinweis{(Präfix: REQ-EXT-COM)}

\section{Produkteigenschaften}
\label{srs:produkteigenschaften}

\hinweis{(Präfix: REQ-PROD)}

\section{Systemattribute}
\label{srs:systemattribute}

\hinweis{(Präfix: REQ-SYS)}

Die Auflistung richtet sich nach ISO-9126 beziehungsweise nach der
neueren Norm ISO/IEC-25012.

\subsection{Änderbarkeit und Wartbarkeit}
\label{srs:aenderbarkeit-wartbarkeit}

\hinweis{(Präfix: REQ-SYS-AW)}

\subsection{Effizienz}
\label{srs:effizienz}

\hinweis{(Präfix: REQ-SYS-EFF)}

\subsection{Übertragbarkeit}
\label{srs:uebertragbarkeit}

\hinweis{(Präfix: REQ-SYS-UE)}

\subsection{Zuverlässigkeit}
\label{srs:zuverlaessigkeit}

\hinweis{(Präfix: REQ-SYS-ZV)}

\subsection{Funktionalität}
\label{srs:funktionalität}

\hinweis{(Präfix: REQ-SYS-FU)}

\subsection{Benutzbarkeit}
\label{srs:benutzbarkeit}

\hinweis{(Präfix: REQ-SYS-USE)}

\section{Datenbankanforderungen}
\label{srs:datenbankanforderungen}

\hinweis{(Präfix: REQ-DB)}

\hinweis{Wenn Anforderungen an eine Datenbank gibt, dann stehen diese hier.}

\chapter{Zusatzmaterial}
\label{srs:zusatzmaterial}

